\section{선별 문제 풀이 II}
\label{sec:symmetric-discriminant-solutions-b}

\subsection*{문제 \thechapter.18: 종결식의 노름 해석}

\begin{solution}
$f$가 일계수 $m$차이므로 나눗셈 알고리즘에 의해 모든 잉여류는
\[
  a_0+a_1x+\cdots+a_{m-1}x^{m-1}
\]
로 유일하게 표현된다. 따라서
\[
  1,x,\dots,x^{m-1}
\]
는 $A=R[X]/(f)$의 $R$-자유기저이다.

실베스터 사상
\[
  \Phi:R[X]_{<n}\oplus R[X]_{<m}	o R[X]_{<m+n},
  \qquad (u,v)\mapsto uf+vg
\]
를 생각하자. 공역의 표준기저 대신
\[
  fX^{n-1},\dots,f,
  X^{m-1},\dots,1
\]
을 사용한다. $f$가 일계수이므로 이 기저에서 표준기저로 가는 행렬은 대각성분이 모두 $1$인 삼각행렬이며 행렬식이 $1$이다.

새 기저에서 $\Phi$는 첫 $n$개의 벡터를 그대로 보내고, 뒤의 $m$개 벡터에는 $g$를 곱한다. $(f)$로 나누면 뒤쪽 블록은 정확히
\[
  \mu_{g(x)}:A\to A,
  \qquad a\mapsto g(x)a
\]
의 행렬이다. 따라서
\[
  \Res(f,g)=\det\Phi
  =1\cdot\det\mu_{g(x)}
  =\operatorname N_{A/R}(g(x)).
\]
몫환에서는 $g$를 $f$로 나눈 나머지만 중요하므로 오른쪽은 $g$의 형식적 차수 선택과 무관하다.
\end{solution}

\subsection*{문제 \thechapter.19: 종결식의 근 공식}

\begin{solution}
먼저 직접 계산하면
\[
  \Res(X-\alpha,g)=g(\alpha).
\]
실제로 $g=b_0X^n+\cdots+b_n$일 때 실베스터 행렬을 일차식의 계수행으로 소거하면 이 값이 남는다.

이제
\[
  f=a\prod_{i=1}^m(X-\alpha_i)
\]
라 하자. 첫 번째 변수에 관한 종결식의 곱셈성과 스케일링 성질을 반복 적용하면
\[
\begin{aligned}
  \Res(f,g)
  &=a^n\prod_{i=1}^m\Res(X-\alpha_i,g)\\
  &=a^n\prod_{i=1}^m g(\alpha_i).
\end{aligned}
\]
또한
\[
  g(\alpha_i)
  =b\prod_{j=1}^n(\alpha_i-\beta_j)
\]
이므로
\[
  \Res(f,g)
  =a^n b^m
  \prod_{i=1}^m\prod_{j=1}^n
  (\alpha_i-\beta_j).
\]
\end{solution}

\subsection*{문제 \thechapter.20: 판별식과 도함수의 종결식}

\begin{solution}
분해환에서
\[
  f(X)=\prod_{i=1}^n(X-\alpha_i)
\]
라 쓰자. 종결식의 근 공식으로
\[
  \Res(f,f')=\prod_{i=1}^n f'(\alpha_i).
\]
곱의 미분법에서
\[
  f'(\alpha_i)=\prod_{j\ne i}(\alpha_i-\alpha_j)
\]
이므로
\[
  \Res(f,f')
  =\prod_{i\ne j}(\alpha_i-\alpha_j).
\]
각 $i<j$에 대해
\[
  (\alpha_i-\alpha_j)(\alpha_j-\alpha_i)
  =-(\alpha_i-\alpha_j)^2
\]
이고 이런 쌍이 $n(n-1)/2$개이므로
\[
  \Res(f,f')
  =(-1)^{n(n-1)/2}
  \prod_{i<j}(\alpha_i-\alpha_j)^2.
\]
오른쪽의 제곱곱은 $\Disc(f)$이다. 따라서
\[
  \Disc(f)=(-1)^{n(n-1)/2}\Res(f,f').
\]
\end{solution}

\subsection*{문제 \thechapter.21: 판별식과 교대군}

\begin{solution}
근들을 $\alpha_1,\dots,\alpha_n$이라 하고
\[
  \delta=\prod_{i<j}(\alpha_i-\alpha_j)
\]
라 두자. 그러면 $\delta^2=\Disc(f)$이다. $\sigma\in G$가 근들에 유도하는 순열을 같은 기호로 쓰면
\[
  \sigma(\delta)=\sgn(\sigma)\delta.
\]
따라서
\[
  G\subseteq A_n
\]
인 것은 모든 $\sigma\in G$가 $\delta$를 고정하는 것과 동치이다. 분해체의 전체 갈루아군의 고정체는 $K$이므로 이는
\[
  \delta\in K
\]
와 동치이다.

$\delta\in K$이면 판별식은 당연히 $K$의 제곱이다. 반대로 $d^2=\Disc(f)$인 $d\in K^\times$가 있다고 하자. $\delta$와 $d$는 $X^2-\Disc(f)$의 근이다. $\charac K\ne2$이므로 이 다항식의 두 근은 $d,-d$이고
\[
  \delta=\pm d\in K.
\]
따라서 세 조건이 동치이다.
\end{solution}

\subsection*{문제 \thechapter.23: 일반 삼차식의 판별식}

\begin{solution}
\[
  f(X)=X^3+aX^2+bX+c
\]
에서
\[
  X=Y-\frac a3
\]
를 대입한다. 전개하면
\[
  f(X)=Y^3+pY+q
\]
이고
\[
  p=b-\frac{a^2}{3},
  \qquad
  q=c-\frac{ab}{3}+rac{2a^3}{27}.
\]
평행이동은 근들의 모든 차이를 보존하므로 판별식도 보존된다. 따라서
\[
  \Disc(f)=-4p^3-27q^2.
\]
분모를 통분하여 전개하면
\[
  \Disc(f)
  =a^2b^2-4b^3-4a^3c-27c^2+18abc.
\]

이 계산은 우선 $\Q(a,b,c)$에서 수행했지만 최종 양변은 $\Z[a,b,c]$의 다항식이다. 대수적으로 독립인 변수에 관한 다항식 항등식이므로, 임의의 가환환에서 $a,b,c$를 어떤 원소로 특수화해도 등식이 유지된다. 따라서 $3$이 가역이 아닌 환에서도 최종 공식은 성립한다.
\end{solution}

\subsection*{문제 \thechapter.28: 두 이차식의 종결식}

\begin{solution}
실베스터 행렬은
\[
  S=
  \begin{pmatrix}
  a_0&a_1&a_2&0\\
  0&a_0&a_1&a_2\\
  b_0&b_1&b_2&0\\
  0&b_0&b_1&b_2
  \end{pmatrix}.
\]
행렬식을 전개하면
\[
\begin{aligned}
  \det S={}&a_0^2b_2^2-a_0a_1b_1b_2
  +a_0a_2b_1^2-2a_0a_2b_0b_2\\
  &+a_1^2b_0b_2-a_1a_2b_0b_1+a_2^2b_0^2.
\end{aligned}
\]
한편
\[
\begin{aligned}
 &(a_0b_2-a_2b_0)^2
 +(a_1b_2-a_2b_1)(a_1b_0-a_0b_1)\\
 ={}&a_0^2b_2^2-2a_0a_2b_0b_2+a_2^2b_0^2\\
 &+a_1^2b_0b_2-a_0a_1b_1b_2
 -a_1a_2b_0b_1+a_0a_2b_1^2,
\end{aligned}
\]
이므로 두 식이 같다. 따라서
\[
  \Res(f,g)
  =(a_0b_2-a_2b_0)^2
  +(a_1b_2-a_2b_1)(a_1b_0-a_0b_1).
\]
\end{solution}

\subsection*{문제 \thechapter.29: 거듭제곱 기저의 판별식}

\begin{solution}
먼저 보편 계수환
\[
  R_0=\Z[s_1,\dots,s_n]
\]
과 다항식
\[
  f_0=X^n-s_1X^{n-1}+\cdots+(-1)^ns_n
\]
을 생각한다. 기본대칭식들이 대수적으로 독립이므로 $R_0$은 다항식환이고 분수체
\[
  K=\Q(s_1,\dots,s_n)
\]
에 들어간다. $f_0$의 분해체에서 근을 $\alpha_1,\dots,\alpha_n$이라 하자.

$A=K[X]/(f_0)$에서 $x$를 $X$의 잉여류라 하면, 분리 단순확장의 대각합 공식으로
\[
  \operatorname{Tr}_{A/K}(x^r)
  =\sum_{k=1}^n\alpha_k^r.
\]
따라서 기저 $1,x,\dots,x^{n-1}$의 대각합 그람행렬은
\[
  G=\left(
  \sum_{k=1}^n\alpha_k^{i+j}
  \right)_{0\le i,j<n}.
\]
반데르몬드 행렬
\[
  V=(\alpha_k^i)_{
  \substack{0\le i<n\\1\le k\le n}}
\]
를 쓰면 $G=VV^{\mathsf T}$이다. 그러므로
\[
\begin{aligned}
  D_{A/K}(1,x,\dots,x^{n-1})
  &=\det G\\
  &=(\det V)^2\\
  &=\prod_{i<j}(\alpha_i-\alpha_j)^2\\
  &=\Disc(f_0).
\end{aligned}
\]
두 변은 계수들의 정수계수 다항식이고 환준동형과 호환된다. 따라서 보편 계수환에서 얻은 항등식을 임의의 가환환 $R$과 일계수다항식 $f$로 특수화하면
\[
  D_{R[X]/(f)/R}(1,x,\dots,x^{n-1})=\Disc(f)
\]
가 성립한다.
\end{solution}
